% Cuerpo íntegro, sin preámbulo ni portada duplicada.



\section*{Objet et résultats}
On démontre que le raffinement réversible avec retenue produit par l’architecture APP--TRIT--TPK admet une forme canonique stable : la somme des formes des cellules raffinées coïncide exactement avec la forme de la cellule initiale et tous les pôles intérieurs s’annulent. L’identité s’étend exactement au complexe produit cuboïdal et conserve les résidus itérés de la frontière externe. Le registre dodécaphasique \(K\) détermine un repère inversible \(O_K\), une métrique de Gram et une archive isométrique qui préservent les volumes extérieurs et permettent de récupérer la trajectoire. Un calcul exhaustif des mineurs montre que \(O_K\) n’est pas totalement non négatif. La direction déjà construite \(u_K=P_3P_{11}K\), combinée à \(\rho_A=\pi_{\rm HMT}/\varphi_{\rm HMT}^2\), définit un flot diagonal positif de déterminant un qui conserve chaque signe de Plücker et chaque strate positroïdale. Cette séparation permet de formuler, sous des conditions réfutables, l’atlas positroïdal intrinsèque du TPK, sa dynamique d’amas et son image directe vers une section amplituédrique. La centralité de l’électron est liée à l’unique point fixe \((5,5)\) de l’atlas APP, sans le confondre avec la donnée nonadique \(111111\) ni avec la signature régionale \(555555\).


\section{Thèse et ordre causal}

L’objet de départ n’est pas une grassmannienne choisie extérieurement. La chaîne
\[
\APP\longrightarrow\mathrm{TRIT}\longrightarrow\TPK
\longrightarrow\widehat x
\longrightarrow\mathcal C_{\rm cont}^{\rm disc}
\]
produit un état enrichi conservant les feuillets, l’orientation, le reste, le quotient, la retenue, le parcours, la frontière, la survie, la mémoire et le retour. Les cinq constructions du continu émergent conjointement. La géométrie positive est une expression de leur raffinement compatible.

L’article distingue trois assertions qui ne doivent pas être fusionnées :
\begin{enumerate}
\item le télescopage et l’archive métrique sont des résultats exacts ;
\item la courbe des moments, \(D_K(s)\) et le flot centré \(g_{K,\rho}(\tau)\) sont de nouvelles formalisations exactes ;
\item l’atlas positroïdal intrinsèque, la graine d’amas et l’égalité avec un amplituèdre complet exigent des applications supplémentaires expressément déclarées.
\end{enumerate}
Cette séparation ne diminue pas le résultat : elle identifie la démonstration de chaque flèche et évite d’attribuer à une matrice métrique une positivité qu’elle ne possède pas.

\section{Théorème directeur}


% BEGIN OWNER /Users/ruben/Documents/New project/output/IMPLEMENTACION_GEOMETRIA_POSITIVA_TPK_20260918/02_TEOREMA_TRONCAL_GEOMETRIA_POSITIVA_TPK.tex
% Cuerpo de inserción autónomo. No modifica ningún manuscrito activo.
\section{Géométrie positive coinductive du raffinement TPK}
\label{geometria_positiva:sec:geom-positiva-tpk}

La structure positive étudiée dans cette section n’est pas introduite à partir d’une grassmannienne préalablement supposée. Elle provient de la même chaîne génératrice qui construit l’état enrichi :
\[
 \APP\longrightarrow\mathrm{TRIT}\longrightarrow\TPK
 \longrightarrow \widehat x
 \longrightarrow\mathcal C_{\rm cont}^{\rm disc}.
\]
APP conserve simultanément ses deux feuillets, ses restes et ses quotients ; TRIT type le régime et oriente chaque transition ; TPK sélectionne, transporte et actualise sans perdre la retenue, le parcours, la frontière, la survie ni la mémoire. Les cinq constructions du continu émergent conjointement de cet état. La réalisation grassmannienne qui suit est un lecteur ultérieur de son raffinement compatible et non une entrée sélectionnant la dynamique.

\subsection{Retenue, subdivision et paramètre ordonné}

Fixons une base entière \(B\ge2\). L’opération élémentaire de raffinement avec retenue \(c\in\{0,\ldots,B-1\}\) est
\begin{equation}
 L_c(x,y)=(Bx-c,By+c).
 \label{geometria_positiva:eq:carry-map-positive}
\end{equation}
Elle conserve la somme remise à l’échelle, puisque
\[
 \frac{Bx-c}{B(x+y)}+\frac{By+c}{B(x+y)}=1.
\]
Son inverse récupère \(c\) par la classe résiduelle déclarée, puis récupère \((x,y)\) par division euclidienne. Par composition, un mot de retenues \(c_1\ldots c_n\) détermine une cellule de profondeur \(n\), et les divisions dans l’ordre inverse restituent le mot complet. La compatibilité entre troncatures produit des cylindres emboîtés et conserve la trajectoire, pas seulement son évaluation limite.

Soit \(M\ge1\) le nombre d’intervalles de la carte initiale et
\begin{equation}
 N_n=B^nM,\qquad t_{n,j}=\frac{j}{N_n},\qquad 0\le j\le N_n.
 \label{geometria_positiva:eq:ordered-carry-points}
\end{equation}
Pour \(0\le j<N_n\), le sous-intervalle d’indice \(c\) du point \(j\) a pour indice
\begin{equation}
 j'=Bj+c,\qquad
 t_{n+1,Bj+c}=t_{n,j}+\frac{c}{BN_n}.
 \label{geometria_positiva:eq:carry-insertion}
\end{equation}
L’inégalité \(j<j'\Rightarrow t_{n,j}<t_{n,j'}\) rend visible l’information nécessaire à la positivité : la retenue ne fait pas que raffiner ; elle insère des points en conservant l’ordre total compatible avec les applications de liaison.

\subsection{Matrice des moments et grassmannienne positive}

\paragraph{Statut.}
La courbe des moments qui suit est une nouvelle formalisation de l’ordre déjà produit par le raffinement ; elle n’est présentée ni comme un opérateur supplémentaire du TPK ni comme une entrée du générateur.

Pour \(1\le r\le N_n+1\), définissons la courbe des moments et sa matrice de colonnes
\[
 v_r(t)=(1,t,\ldots,t^{r-1})^{\mathsf T},\qquad
 Z_n^{(r)}=(v_r(t_{n,0})\ \cdots\ v_r(t_{n,N_n})).
\]

\begin{theorem}[Positivité de Plücker engendrée par la retenue]
\label{geometria_positiva:thm:tpk-vandermonde-positive}
Pour tout \(I=\{j_1<\cdots<j_r\}\),
\begin{equation}
 \Delta_I(Z_n^{(r)})
 =\prod_{1\le a<b\le r}(t_{n,j_b}-t_{n,j_a})>0.
 \label{geometria_positiva:eq:vandermonde-minor}
\end{equation}
La classe des lignes de \(Z_n^{(r)}\) appartient donc à \(\operatorname{Gr}_{>0}(r,N_n+1)\). Les insertions \eqref{geometria_positiva:eq:carry-insertion} prolongent ce système sans modifier le signe d’aucun mineur antécédent.
\end{theorem}
\begin{proof}
Le mineur est un déterminant de Vandermonde. Chaque facteur de \eqref{geometria_positiva:eq:vandermonde-minor} est positif par l’ordre strict des points. Une insertion maintient l’ordre des anciennes colonnes et ajoute des points dans leurs intervalles. Les mineurs formés uniquement de colonnes antécédentes conservent exactement leur valeur ; ceux contenant une nouvelle colonne sont de nouveau des produits de différences positives.
\end{proof}

La version pondérée conserve la même orientation. Si \(\lambda_{n,j}>0\), posons
\[
 C_{n;aj}=\lambda_{n,j}t_{n,j}^{a-1}.
\]
Alors
\begin{equation}
 \Delta_I(C_n)=\left(\prod_{j\in I}\lambda_{n,j}\right)
 \Delta_I(Z_n^{(r)})>0.
 \label{geometria_positiva:eq:weighted-vandermonde}
\end{equation}
Les poids peuvent recevoir une mesure cylindrique compatible ou la racine d’une multiplicité de parcours, à condition d’être déclarés avant l’évaluation de l’objet externe.

\subsection{Masques de survie et réalisation positroïdale}

Le TPK conserve à chaque profondeur un masque de survie \(S_n\subseteq\{0,\ldots,N_n\}\). Pour le rendre visible sans changer le générateur, on autorise
\[
 \lambda_{n,j}>0\quad(j\in S_n),\qquad
 \lambda_{n,j}=0\quad(j\notin S_n).
\]

\paragraph{Statut.}
L’assignation d’une matrice des moments à un masque de survie est une \emph{nouvelle formalisation} : elle utilise des données déjà engendrées par le TPK, mais n’est pas présentée comme un théorème préalablement contenu dans le corpus. Son contenu mathématique exact est le suivant.

\begin{proposition}[Réalisation positroïdale d’un masque]
\label{geometria_positiva:thm:tpk-positroid-cell} Supposons \(|S_n|\ge r\). La matrice pondérée satisfait
\[
 \Delta_I(C_{n,S_n})>0\quad\Longleftrightarrow\quad I\subseteq S_n,
 \qquad
 \Delta_I(C_{n,S_n})=0\quad\Longleftrightarrow\quad I\not\subseteq S_n.
\]
Son matroïde est le matroïde uniforme de rang \(r\) sur \(S_n\), avec les positions complémentaires comme boucles ; en particulier, c’est un positroïde. Les réalisations de deux niveaux sont compatibles dès que leurs masques le sont sous l’application de liaison et que les anciennes colonnes conservent leur identification.
\end{proposition}
\begin{proof}
La factorisation \eqref{geometria_positiva:eq:weighted-vandermonde} montre qu’un mineur s’annule exactement lorsqu’il contient un poids nul. Sur \(S_n\), tous les facteurs sont positifs et tous les sous-ensembles de cardinal \(r\) sont des bases. Un matroïde uniforme réalisé par des points ordonnés sur la courbe des moments est positroïdal ; l’ajout de boucles conserve cette classe. La compatibilité vient de ce qu’un masque raffiné hérite des survivants et ajoute uniquement les colonnes permises par l’hypothèse de compatibilité.
\end{proof}

La proposition certifie la réalisation de chaque masque dans le foncteur des moments ordonnés. Elle ne certifie pas encore un atlas positroïdal intrinsèque du TPK ni la couverture de toutes les cellules de \(\operatorname{Gr}_{\ge0}(r,N_n+1)\). Cette promotion exige un réseau planaire ou une mesure de frontière dérivée du TPK, avec sa permutation décorée ou son collier de Grassmann, et la naturalité de ces données sous raffinement. L’absence de l’une de ces composantes réfuterait la promotion globale, non la proposition précédente.

\subsection{Échanges de Plücker et cartes d’amas}

Pour \(r=2\), écrivons \(z_j=\lambda_j(1,t_j)^{\mathsf T}\). Si \(a<b<c<d\), on a
\begin{equation}
 \Delta_{ac}\Delta_{bd}
 =\Delta_{ab}\Delta_{cd}+\Delta_{ad}\Delta_{bc}.
 \label{geometria_positiva:eq:ptolemy-tpk}
\end{equation}

\paragraph{Statut.}
La relation suivante est exacte pour la réalisation des moments ordonnés. Sa lecture comme structure d’amas du TPK est une \emph{nouvelle formalisation} ; elle n’identifie pas la retenue à une mutation et n’affirme pas que le corpus ait déjà construit un graphe plabique universel.

\begin{proposition}[Échange positif en rang deux]
\label{geometria_positiva:thm:tpk-cluster-rank2} Toutes les coordonnées de \eqref{geometria_positiva:eq:ptolemy-tpk} sont positives et la mutation isolant l’une des deux diagonales est sans soustraction. L’insertion d’un point raffiné par \eqref{geometria_positiva:eq:carry-insertion} commute avec l’évaluation de la relation de Ptolémée. Les triangulations du polygone ordonné et leurs basculements fournissent donc des cartes d’amas positives pour la réalisation ordonnée de chaque niveau.
\end{proposition}
\begin{proof}
L’identité s’obtient en substituant \(\Delta_{ij}=\lambda_i\lambda_j(t_j-t_i)\) et en développant. Tous ses termes sont positifs. L’insertion conserve les anciennes coordonnées et assigne à la nouvelle colonne le paramètre ordonné exact ; évaluer d’abord le déterminant ou d’abord la relation d’échange produit la même identité polynomiale.
\end{proof}

En rang général, pour \(|I|=r-2\) et \(a<b<c<d\) hors de \(I\),
\begin{equation}
 \Delta_{Iac}\Delta_{Ibd}
 =\Delta_{Iab}\Delta_{Icd}+\Delta_{Iad}\Delta_{Ibc}.
 \label{geometria_positiva:eq:short-plucker-tpk}
\end{equation}
Tous les termes sont positifs dans la cellule supérieure et non négatifs dans les cellules de survie. Ainsi, chaque échange court de Plücker est compatible avec le raffinement dans cette réalisation. L’assertion correcte est que la retenue insère des données et que les mutations changent de carte sur ces données ; l’étape de retenue n’est pas identifiée à un basculement. Une compatibilité d’amas intrinsèque au TPK exige en outre de dériver une graine, son carquois ou son graphe plabique et l’action des applications de liaison sur les mutations.

\subsection{Formes canoniques et annulation des frontières intérieures}

Pour un intervalle orienté \(I=[a,b]\), sa forme canonique est
\begin{equation}
 \Omega_I=\frac{dt}{t-a}-\frac{dt}{t-b}.
 \label{geometria_positiva:eq:interval-canonical-form}
\end{equation}
Si \(a=c_0<c_1<\cdots<c_m=b\), alors
\begin{equation}
 \sum_{q=0}^{m-1}\Omega_{[c_q,c_{q+1}]}=\Omega_{[a,b]}.
 \label{geometria_positiva:eq:interval-telescopy}
\end{equation}
Chaque pôle intérieur apparaît deux fois, avec des résidus opposés.

L’identité unidimensionnelle est exacte et ne dépend pas d’une interprétation positroïdale supplémentaire. Son extension cubique est formulée ci-dessous comme une \emph{nouvelle formalisation exacte}. Soit \(Q=\prod_{\ell=1}^d[a_\ell,b_\ell]\) d’orientation \(dx_1\wedge\cdots\wedge dx_d\). Définissons
\begin{equation}
 \Omega_Q=\bigwedge_{\ell=1}^d
 \left(\frac{dx_\ell}{x_\ell-a_\ell}
       -\frac{dx_\ell}{x_\ell-b_\ell}\right).
 \label{geometria_positiva:eq:box-canonical-form}
\end{equation}

\begin{theorem}[Nouvelle formalisation : télescopage cubique multidimensionnel]
\label{geometria_positiva:thm:tpk-higher-residues} Subdivisons chaque intervalle de coordonnées par les points du raffinement TPK et notons \(Q_{\mathbf j}\) les cuboïdes obtenus, munis de l’orientation induite. Alors
\begin{equation}
 \boxed{\displaystyle
 \sum_{\mathbf j}\Omega_{Q_{\mathbf j}}=\Omega_Q.}
 \label{geometria_positiva:eq:box-telescopy}
\end{equation}
Tous les pôles des facettes intérieures s’annulent. Si \(F_1\supset F_2\supset\cdots\supset F_s\) est un drapeau de faces orientées, alors
\begin{equation}
 \operatorname{Res}_{F_s}\cdots\operatorname{Res}_{F_1}\Omega_Q
 =\operatorname{sgn}(F_\bullet)\,\Omega_{F_s}.
 \label{geometria_positiva:eq:iterated-residue}
\end{equation}
En particulier, le résidu à chaque sommet est de module un et de signe fixé par l’orientation.
\end{theorem}
\begin{proof}
La somme sur la grille est le produit extérieur des \(d\) sommes unidimensionnelles. Appliquer \eqref{geometria_positiva:eq:interval-telescopy} dans chaque facteur donne \eqref{geometria_positiva:eq:box-telescopy}. Une facette intérieure possède deux cellules incidentes et des orientations normales opposées ; leurs résidus s’annulent donc. Prendre un résidu élimine le facteur normal correspondant et laisse la forme canonique de la face, avec le signe d’orientation. L’itération démontre \eqref{geometria_positiva:eq:iterated-residue}.
\end{proof}

\begin{corollary}[Indépendance du raffinement]
La forme globale ne dépend pas du nombre de niveaux de retenue insérés. Raffiner change la triangulation ou la cubulation et préserve la somme des formes. La mémoire TPK conserve les cellules intermédiaires traversées, même si la forme globale annule leurs frontières internes.
\end{corollary}

\begin{proposition}[Critère orienté général d’annulation]
\label{geometria_positiva:prop:oriented-residue-chain} Soit \(X=\bigcup_\sigma X_\sigma\) une subdivision positive pure et orientée. Supposons que chaque cellule possède une forme \(\Omega_\sigma\) et que, pour toute facette \(F\subset\overline{X_\sigma}\),
\[
\operatorname{Res}_F\Omega_\sigma=[\sigma:F]\Omega_F.
\]
Pour une chaîne \(c=\sum_\sigma w_\sigma[\sigma]\),
\begin{equation}
\operatorname{Res}_F\!\left(\sum_\sigma w_\sigma\Omega_\sigma\right)
=(\partial c)_F\Omega_F.
\label{geometria_positiva:eq:residue-chain-map}
\end{equation}
En particulier, une facette intérieure s’annule si et seulement si son coefficient dans \(\partial c\) est nul. Les résidus itérés sur les drapeaux intérieurs s’annulent par \(\partial^2=0\).
\end{proposition}
\begin{proof}
La linéarité du résidu donne
\[
\operatorname{Res}_F\!\left(\sum_\sigma w_\sigma\Omega_\sigma\right)
=\sum_{\sigma\supset F}w_\sigma[\sigma:F]\Omega_F,
\]
et le coefficient de \(\Omega_F\) est précisément celui de \(F\) dans \(\partial c\). L’itération identifie la somme des incidences de codimension deux au coefficient correspondant de \(\partial^2c\), qui est nul.
\end{proof}

La proposition est une nouvelle formalisation exacte. Pour identifier la somme à l’unique forme canonique de \(X\), il faut en outre exclure une forme de degré maximal régulière ambiguë. Le corpus fournit déjà le complexe de chemins et de mémoire avec \(D^2=0\) ; la promotion positroïdale exige de construire le morphisme
\[
\Omega_{\rm TPK}:C_\bullet^{\rm TPK}\longrightarrow
\Omega^\bullet_{\log},
\qquad
\operatorname{Res}\circ\Omega_{\rm TPK}
=\Omega_{\rm TPK}\circ D.
\]
Sur les cylindres et les produits, ce morphisme est donné par \eqref{geometria_positiva:eq:interval-canonical-form} et \eqref{geometria_positiva:eq:box-canonical-form}. Pour un atlas positroïdal général, il faut ajouter les strates, incidences, orientations, l’unicité et la compatibilité Plücker--amas. Les poids probabilistes des trajectoires ne peuvent pas non plus multiplier indistinctement les formes : une facette partagée conserverait un résidu proportionnel à la différence des poids. Les facteurs \(\sqrt\mu\) d’une matrice positive relèvent d’une autre opération.

\subsection{Le registre \(K\) comme repère métrique et archivistique}

Le registre dodécaphasique est produit par
\[
 x_{\rm term}\mapsto\mathcal R_{12}
 \mapsto(b^{90},b^{120},Q_{\rm TPK})
 \mapsto U_{\rm sgn}\mapsto K
\]
et vaut
\[
 K=(234,543,140,729,659,824,621,58,914,794,146,601).
\]
Soit \(S\) le décalage cyclique de douze positions et
\[
 O_K=(K,SK,\ldots,S^{11}K),\qquad G_K=O_K^*O_K.
\]
Le développement de référence démontre
\begin{equation}
 \det O_K=5483119259214020183850397930008283625>0,
 \qquad
 589609I\le G_K\le6263^2I.
 \label{geometria_positiva:eq:K-frame-bounds}
\end{equation}
En particulier, \(O_K\) est inversible et définit un repère quantitativement contrôlé. Son déterminant étant positif, sa décomposition polaire peut s’écrire
\[
 O_K=F_KP_K,\qquad P_K=G_K^{1/2}>0,\qquad F_K\in SO(12).
\]

\begin{proposition}[Covariance exacte des données de Plücker]
\label{geometria_positiva:thm:K-plucker-transport} Soit \(X\in\operatorname{Mat}_{r\times12}(\mathbb R)\) de rang \(r\) et \(Y=XO_K^{\mathsf T}\). Alors
\begin{equation}
 \Delta_I(Y)=
 \sum_{\substack{J\subseteq\{1,\ldots,12\}\\ |J|=r}}
 \Delta_J(X)\det\!\bigl((O_K)_{I,J}\bigr).
 \label{geometria_positiva:eq:K-cauchy-binet}
\end{equation}
De manière équivalente, le vecteur de Plücker se transforme par \(\bigwedge^rO_K\). En rang douze,
\[
 \Delta_{\{1,\ldots,12\}}(XO_K^{\mathsf T})
 =(\det O_K)\Delta_{\{1,\ldots,12\}}(X).
\]
\end{proposition}
\begin{proof}
La formule \eqref{geometria_positiva:eq:K-cauchy-binet} est l’identité de Cauchy--Binet appliquée à chaque ensemble de colonnes de \(Y\). La formulation extérieure et le cas de rang douze en sont respectivement les expressions fonctorielle et maximale.
\end{proof}

La covariance de Plücker n’équivaut pas à la préservation de la positivité totale standard. En effet, \(O_K\) possède des mineurs des deux signes. Deux témoins explicites sont
\begin{equation}
 \det\begin{pmatrix}234&543\\543&140\end{pmatrix}=-262089,
 \qquad
 \det\begin{pmatrix}234&140\\543&729\end{pmatrix}=94566.
 \label{geometria_positiva:eq:K-mixed-minors}
\end{equation}
Par exemple, l’image par \(O_K\) du plan de coordonnées positif engendré par les deux premiers vecteurs canoniques contient déjà le premier mineur négatif de \eqref{geometria_positiva:eq:K-mixed-minors}. Par conséquent,
\begin{equation}
 \left\{[XO_K^{\mathsf T}]:
 [X]\in\operatorname{Gr}_{\ge0}(r,12)\right\}
 \not\subseteq\operatorname{Gr}_{\ge0}(r,12)
 \quad\text{en général}.
 \label{geometria_positiva:eq:K-not-standard-totally-positive}
\end{equation}
Ainsi, \(O_K\) est un repère métrique et archivistique inversible ; ce n’est ni une matrice de mesure de frontière ni un générateur de positivité totale standard.

\paragraph{Carte positive adaptée à \(K\).}
L’image exacte
\begin{equation}
 \operatorname{Gr}_{\ge0}^{K}(r,12)
 :=\left\{[XO_K^{\mathsf T}]:
 [X]\in\operatorname{Gr}_{\ge0}(r,12)\right\}
 \label{geometria_positiva:eq:K-adapted-positive-chart}
\end{equation}
définit une carte positive adaptée au repère. Si \(\Phi_K([X])=[XO_K^{\mathsf T}]\), ses coordonnées adaptées sont
\begin{equation}
 \Delta_I^K([Y])
 :=\Delta_I\!\left(Y(O_K^{\mathsf T})^{-1}\right)\ge0.
 \label{geometria_positiva:eq:K-adapted-plucker}
\end{equation}
Cette positivité est exacte par transport de structure ; elle ne doit pas être confondue avec la positivité des mineurs de \(Y\) dans la base ambiante standard. De même, une forme canonique \(\Omega\) dans la carte d’origine se transporte comme \((\Phi_K)_*\Omega\).

\paragraph{Nouvelle formalisation : action diagonale positive de \(K\).}
Comme toutes les composantes \(K_i\) sont strictement positives, on définit
\begin{equation}
 D_K(s):=\operatorname{diag}
 \bigl(e^{s\log K_1},\ldots,e^{s\log K_{12}}\bigr),
 \qquad s\in\mathbb R,
 \label{geometria_positiva:eq:K-positive-diagonal-flow}
\end{equation}
de sorte que \(D_K(1)=\operatorname{diag}(K_1,\ldots,K_{12})\). Pour toute matrice \(X\in\operatorname{Mat}_{r\times12}\),
\begin{equation}
 \Delta_I(XD_K(s))
 =e^{\,s\sum_{i\in I}\log K_i}\Delta_I(X).
 \label{geometria_positiva:eq:K-diagonal-minors}
\end{equation}
Cette action préserve exactement \(\operatorname{Gr}_{\ge0}(r,12)\) et sa stratification par annulations de mineurs. L’action diagonale est une nouvelle formalisation construite à partir du registre engendré \(K\) ; elle n’est ni attribuée au développement original de référence de \(O_K\) ni utilisée pour masquer les signes mixtes de \eqref{geometria_positiva:eq:K-mixed-minors}.

\paragraph{Nouvelle formalisation : flot centré déterminé par \(K\).}
Le corpus construit
\begin{equation}
 u_K=P_3P_{11}K,\qquad
 \|u_K\|^2=
 \frac{6638585+2275584\sqrt5}{20}>0,
 \qquad u_K\perp\mathbf1.
 \label{geometria_positiva:eq:K-centered-direction}
\end{equation}
Soit \(\widehat u_K=u_K/\|u_K\|\) et \(\rho_A=\pi_{\rm HMT}/\varphi_{\rm HMT}^2>0\). Alors
\begin{equation}
 g_{K,\rho}(\tau)=
 \operatorname{diag}\!\left(
 \rho_A^{\tau\widehat u_{K,1}},\ldots,
 \rho_A^{\tau\widehat u_{K,12}}\right)
 \in SL(12,\mathbb R)_{>0}.
 \label{geometria_positiva:eq:K-centered-positive-flow}
\end{equation}
Pour toute matrice \(X\),
\begin{equation}
 \Delta_I(Xg_{K,\rho}(\tau))
 =\rho_A^{\tau\sum_{i\in I}\widehat u_{K,i}}\Delta_I(X).
 \label{geometria_positiva:eq:K-centered-flow-minors}
\end{equation}
Le flot préserve donc les signes et les mineurs nuls. Le registre \(K\) fixe la direction et le rapport engendré \(\rho_A\) normalise la vitesse. Le rapport \(\rho_A\) est évalué ici à une étape ultérieure : \(\pi_{\rm HMT}\) et \(\varphi_{\rm HMT}\) sont deux coordonnées de la quadrirelation \((\pi,\varphi,e,\alpha)\) engendrée coinductivement par le même état APP--TRIT--TPK\@. Ni leurs valeurs conventionnelles ni le rapport déjà évalué ne sélectionnent \(K\), \(u_K\) ou le raffinement. L’identification de \(\tau\) au temps physique exige une horloge \(\tau(t)\) ; le prolongement à toute profondeur exige un cocycle sectoriel compatible avec \(j\mapsto Bj+c\).

L’archive complète
\[
 \mathcal Vv=\left(A_\infty^{1/2}v,(D_jR_jv)_{j\ge0}\right),
 \qquad\mathcal V^*\mathcal V=I,
\]
conserve cette information à tous les degrés :
\begin{equation}
 \left\langle\bigwedge^r\mathcal V\,\xi,
 \bigwedge^r\mathcal V\,\eta\right\rangle
 =\langle\xi,\eta\rangle,
 \qquad
 \|\wedge^r\mathcal V\,\xi\|=\|\xi\|.
 \label{geometria_positiva:eq:exterior-isometry}
\end{equation}
Elle conserve donc les déterminants de Gram et les normes de Plücker. Cette isométrie n’affirme pas que \(\mathcal V\) transforme les mineurs de \(O_K\) en mineurs positifs : elle préserve l’information archivée dans l’espace et l’orientation explicitement transportés.

La lecture rationnelle
\[
 \kappa_{\rm ag}=
 \frac{234543140729659824621058914794146601}{10^{36}-1}
\]
est réversible une fois fixées la base, la longueur et l’origine de phase. La fonction de couplage de \(K\) s’exprime donc en trois couches distinctes : sa lecture arithmétique \(\kappa_{\rm ag}\), sa forme de Gram positive \(G_K\) et son transport d’incidence. Toutes trois sont des expressions du registre engendré ; aucune n’est utilisée pour choisir rétrospectivement la trajectoire.

\subsection{Flot aréal--volumétrique et rapport marqué}

Le calendrier TPK induit la matrice primitive
\[
 F_{\rm av}=\begin{pmatrix}0&1\\1&1\end{pmatrix},
\]
dont les puissances obéissent à la récurrence de Fibonacci. La mesure de Parry de son enveloppe symbolique distingue les modes aréal et volumétrique et satisfait
\begin{equation}
 \frac{\mu_{\rm ar}}{\mu_{\rm vol}}=\varphi_{\rm HMT}^{-2}.
 \label{geometria_positiva:eq:area-volume-ratio}
\end{equation}
La lecture régionale marque une seule fois ce rapport :
\begin{equation}
 \Theta_n=\pi_{\rm HMT}\frac{F_{n-1}}{F_{n+1}}
 \longrightarrow\frac{\pi_{\rm HMT}}{\varphi_{\rm HMT}^2}.
 \label{geometria_positiva:eq:marked-area-volume}
\end{equation}
Le lecteur réciproque \(R(x,y)=x/y^2\), conjugué par \(J(x,y)=(y,x)\), produit l’orientation \(\varphi_{\rm HMT}/\pi_{\rm HMT}^2\). Ces rapports sont des lecteurs typés du flot ; ils ne sont pas identifiés à \(\kappa_{\rm ag}\), bien qu’ils participent tous à la même architecture de couplage.

La forme canonique du complexe est stable sous raffinement, tandis que \eqref{geometria_positiva:eq:area-volume-ratio} détermine la proportion stationnaire de ses deux modes. Le cercle se referme ainsi : le raffinement arithmétique engendre la subdivision ; la raison holographique gouverne sa distribution asymptotique ; \(K\) conserve la métrique et l’orientation ; et la forme canonique ne retient que la frontière externe après annulation des parois internes.

\subsection{Centre électronique et section matérielle}

Normalisons l’atlas \(9\times9\) par
\[
 \nu(r,s)=\left(\frac{r-1}{8},\frac{s-1}{8}\right).
\]
L’inversion cellulaire \(\rho_c(r,s)=(10-r,10-s)\) devient \((u,v)\mapsto(1-u,1-v)\). Son unique point fixe est
\begin{equation}
 \nu(5,5)=\left(\frac12,\frac12\right).
 \label{geometria_positiva:eq:electron-center}
\end{equation}
Les \(81\) états constituent des points d’un maillage dont la cubulation élémentaire possède \(64\) rectangles ; ils ne sont pas présentés comme \(81\) cellules bidimensionnelles.

\begin{theorem}[Unicité de la section ponctuelle équivariante]
\label{geometria_positiva:thm:electron-equivariant-section} Dans une fibre stable sur laquelle l’étiquette grossière demeure fixe, toute sélection ponctuelle équivariante pour \(\rho_c\) doit prendre la valeur \((5,5)\). La fibre électronique de niveau zéro contient cet unique point ; aucune autre fibre n’admet une telle sélection.
\end{theorem}
\begin{proof}
Si \(a(y)\) est équivariante et \(y\) est fixe, alors \(a(y)=a(\rho_cy)=\rho_ca(y)\). L’unique solution de \(\rho_c(r,s)=(r,s)\) est \(r=s=5\). Le dénombrement de l’atlas situe cette cellule dans la fibre électronique de niveau zéro.
\end{proof}

Le théorème explique pourquoi l’électron est une section ponctuelle centrale et les autres espèces sont présentées par des orbites ou des multisections. En particulier, les objets typés suivants doivent rester distincts :
\begin{enumerate}
 \item la cellule fixe \((5,5)\), appartenant à l’atlas \(9\times9\) ;
 \item la donnée \(111111\), appartenant à la cinquième frontière nonadique ;
 \item la signature multiplicative transversale \(\mathcal E_{\times}=555555\), appartenant au lecteur régional de \(\pi\) ;
 \item le bloc de signe \((+1,+1,+1)\), qui apparaît dans le mot dodécaphasique
       \[
        \tau_e=(+1,+1,+1,-1,-1,-1,+1,+1,+1,-1,-1,-1).
       \]
\end{enumerate}
La structure APP centrale, distincte de ces trois registres, est
\[
 \mathcal A(5,5)=((1,1),(7,2)),\qquad
 5+5=1+9\cdot1,\qquad 5\cdot5=7+9\cdot2.
\]
La coïncidence ou la répétition de chiffres ne constitue pas une identification : \((5,5)\), \(111111\) et \(555555\) possèdent des domaines, opérations et fonctions causales différents.

La positivité matérielle est ultérieure. Si \(M_M^{(0)}\ge0\) et \(R_{\rm act}>0\), alors
\begin{equation}
 M_M^{\beta,r}=R_{\rm act}^{B_M^r/2}
 M_M^{(0)}R_{\rm act}^{B_M^r/2}\ge0.
 \label{geometria_positiva:eq:material-positive-congruence}
\end{equation}
La congruence conserve le cône positif et réalise sur la branche de masse la positivité précédemment structurée ; elle ne transforme pas la masse observée en entrée de l’atlas.

\subsection{Multimorphologie positive de l’état enrichi}

La multimorphologie ne désigne pas une superposition informelle de figures. C’est un système de morphismes typés de source commune :
\[
\begin{aligned}
\mathfrak P_n&:\widehat{\mathcal X}_{\rm TPK,n}
\longrightarrow\operatorname{Gr}_{\ge0}(k,N_n),\\
\mathfrak R_n&:\widehat{\mathcal X}_{\rm TPK,n}
\longrightarrow C^+_{d,n},\\
\mathfrak C_n&:\widehat{\mathcal X}_{\rm TPK,n}
\longrightarrow\mathcal Q^+_{\rm con,n},\\
\mathfrak M_n&:\widehat{\mathcal X}_{\rm TPK,n}
\longrightarrow\mathcal L^+_{M,n}.
\end{aligned}
\]
Leurs images sont respectivement les positivités de Plücker, réticulaire, constitutive et opératoire de masse. Chaque morphisme conserve les coordonnées nécessaires à son lecteur et commute avec la liaison de raffinement de son codomaine. Les quatre codomaines ne sont ni identifiés ni utilisés pour se sélectionner mutuellement.

\begin{proposition}[Cohérence multimorphologique]
\label{geometria_positiva:prop:multimorphology-coherence} Si chaque morphisme commute avec la troncature, la frontière et le résidu, l’annulation des faces intérieures se transporte aux quatre réalisations. Si, de plus, une application cinématique
\[
\mathfrak A_{Z,n}:C\longmapsto CZ^{\mathsf T}
\]
est de degré un, préserve l’orientation et couvre son image avec des frontières compatibles, la forme amplituédrique induite est indépendante de la profondeur de représentation.
\end{proposition}
\begin{proof}
La proposition~\ref{geometria_positiva:prop:oriented-residue-chain} s’applique avant chaque réalisation. Les carrés de compatibilité transportent le coefficient d’incidence et le résidu. La condition de degré un évite les multiplicités intérieures dans l’image directe cinématique.
\end{proof}

La réalisation \(\mathfrak P_n\) par courbe des moments est exacte comme formalisation ordonnée ; sa promotion en atlas positroïdal intrinsèque conserve les conditions déjà déclarées. Les branches réticulaire et opératoire possèdent des développements HMT--MD de référence antérieurs. La branche constitutive doit préciser, pour chaque application, sa forme positive et son lecteur. La multimorphologie réunit ces applications ; elle n’ajoute pas une cinquième positivité.

% FORMALIZACION_REUNIDA: composición de la respuesta constitutiva de III
% con el lector positivo de GEOMETRIA_POSITIVA/08_ARTICULO_AUTONOMO.tex.
% Propietario de las operaciones: III/manuscrito__sections__04_respuesta_constitutiva.tex.
% Residencia propuesta: después de la coherencia multimorfológica, antes del
% empuje amplituhedral. No modifica sus hipótesis de frontera, residuo o grado.
\subsection{Réalisation constitutive positive et naturalité spectrale}
\label{empalme:positividad-constitutiva}

La réalisation constitutive de la proposition~\ref{geometria_positiva:prop:multimorphology-coherence} reçoit un opérateur déjà construit à partir des deux feuillets orientés du TPK. La section~\ref{iii:iii:sec:constitutiva} détermine sa réponse et son inverse ; la présente composition explicite la forme positive et la liaison de raffinement utilisées par cette réalisation.

Soit $\widehat{\mathcal X}_{\rm con}$ le domaine des états enrichis dont la lecture angulaire appartient à la chambre $x>y>0$. Pour $\xi\in\widehat{\mathcal X}_{\rm con}$, la lecture produit un espace à deux feuillets $E_\xi$, une involution orthogonale $\mathcal R_\xi$ et le transport
\[
 P_{\pm,\xi}=\frac{I\pm\mathcal R_\xi}{2},\qquad
 T_\xi=q_{+,\xi}P_{+,\xi}+q_{-,\xi}P_{-,\xi},
 \qquad 0<q_{+,\xi}<q_{-,\xi}<1.
\]
Les incidences $90$ et $120$ proviennent du calendrier et agissent sur ce même transport. La réponse est
\[
 V_\xi=(I-T_\xi^{90})(I-T_\xi^{120})^{-1},\qquad
 W_{m,\xi}=V_\xi\otimes I_{9^m}.
\]
L’amplification conserve toutes les positions de la nouvelle fibre nonadique. Sur $E_{m,\xi}=E_\xi\otimes\mathbb C^{9^m}$, on définit la forme
\begin{equation}
 \mathfrak C_m(\xi)(v,w)
 :=\langle W_{m,\xi}v,W_{m,\xi}w\rangle
 =\langle v,W_{m,\xi}^{2}w\rangle.
 \label{empalme:eq:forma-constitutiva}
\end{equation}
Il s’agit d’une réalisation concrète dans le cône des formes hermitiennes définies positives. Le carré conserve les réponses quadratiques des deux feuillets exprimant les coordonnées $\widehat\varepsilon$ et $\widehat\mu$ de \eqref{iii:iii:eq:four}.

\begin{proposition}[Coercivité, récupération et liaison nonadique]
\label{empalme:prop:constitutiva-natural} La forme~\eqref{empalme:eq:forma-constitutiva} satisfait, pour tout $v\in E_{m,\xi}$ non nul,
\[
 \frac9{16}\|v\|^2
 <\mathfrak C_m(\xi)(v,v)<\|v\|^2.
\]
Si $A_{m,\xi}=W_{m,\xi}^{2}$ et $\mathbb E_m$ est l’espérance partielle normalisée sur le dernier facteur de dimension neuf, alors
\[
 \mathbb E_m(A_{m+1,\xi})=A_{m,\xi}.
\]
La forme, avec les étiquettes orientées des feuillets, récupère $V_\xi$, les deux canaux $q_{\pm,\xi}$ et leurs coordonnées angulaires. Si un transport unitaire $J:E_\xi\to E_{\xi'}$ satisfait $JT_\xi=T_{\xi'}J$, alors
\[
 (J\otimes I)^*A_{m,\xi'}(J\otimes I)=A_{m,\xi}.
\]
L’identité demeure valable pour toute fonction continue du spectre appliquée à la réponse avant amplification.
\end{proposition}

\begin{proof}
Les projecteurs de feuillet sont orthogonaux et leur somme est l’identité. Le calcul spectral de~\eqref{iii:iii:eq:rchannels} donne
\[
 V_\xi=r_{+,\xi}P_{+,\xi}+r_{-,\xi}P_{-,\xi},\qquad
 r_{\pm,\xi}=f(q_{\pm,\xi}^{30}),\qquad
 f(s)=\frac{1+s+s^2}{1+s+s^2+s^3}.
\]
Dans $0<s<1$, on a
\[
 4(1+s+s^2)-3(1+s+s^2+s^3)
 =(1-s)(3s^2+2s+1)>0,
\]
et le dénominateur dépasse le numérateur de $s^3>0$. Donc $3/4<r_{\pm,\xi}<1$. La décomposition orthogonale et le carré de ces bornes démontrent la coercivité, qui demeure identique sous amplification.

Comme $A_{m+1,\xi}=A_{m,\xi}\otimes I_9$ et que la trace normalisée du dernier facteur satisfait $\operatorname{tr}_9(I_9)=1$,
\[
 \mathbb E_m(A_{m+1,\xi})
 =(\operatorname{id}\otimes\operatorname{tr}_9)
 (A_{m,\xi}\otimes I_9)=A_{m,\xi}.
\]
La même égalité vaut pour chaque polynôme en $W_{m,\xi}$. L’approximation uniforme sur le spectre compact et la continuité de l’espérance l’étendent à toute fonction continue.

La racine positive de $A_{0,\xi}$ récupère $V_\xi$. Les projecteurs orientés identifient ses deux valeurs propres. La fonction $f$ est strictement décroissante dans $(0,1)$ par \eqref{iii:iii:eq:monotone} ; son inverse détermine les racines uniques de~\eqref{iii:iii:eq:cubic}. La racine positive d’ordre trente restitue $q_{\pm,\xi}$, et \eqref{iii:iii:eq:recover-angular} restitue $x$ et $y$.

Enfin, $JT_\xi=T_{\xi'}J$ implique la même identité pour toutes ses puissances. L’inversibilité de $I-T^{120}$ fournit
\[
 J(I-T_\xi^{120})^{-1}
 =(I-T_{\xi'}^{120})^{-1}J.
\]
Il en résulte $JV_\xi=V_{\xi'}J$. L’amplification, le carré et l’unitarité de $J$ donnent la congruence déclarée. L’argument polynomial et son extension uniforme démontrent la dernière assertion.
\end{proof}

La liaison de cette construction est l’espérance normalisée conservant l’opérateur amplifié. La compatibilité avec un transport géométrique utilise l’entrelacement $JT_\xi=T_{\xi'}J$ exprimé dans l’énoncé. Pour sa part, le transport des formes canoniques et des résidus maintient les carrés de frontière de la proposition~\ref{geometria_positiva:prop:multimorphology-coherence}. Les opérations de chaque niveau sont ainsi précisées : réponse angulaire, forme constitutive positive, amplification de la fibre et transport des incidences. La restitution angulaire récupère cette lecture ; l’archive enrichie conserve en outre la trajectoire qui l’a produite.

\subsection{Formalisation amplituédrique conditionnelle}

\paragraph{Statut.}
Ce qui suit est une \emph{nouvelle formalisation conditionnelle}. Le corpus fournit le raffinement, les parcours, la survie et la donnée ordonnée ; l’identification de son image à un amplituèdre et la promotion de la somme cellulaire en forme canonique exigent les hypothèses géométriques explicitement énoncées.

Fixons \(k,m\ge1\) et \(N=N_n+1\ge k+m\). Soit \(Z_n=Z_n^{(k+m)}\) la donnée externe positive construite par la retenue et soit \(C_n\in\operatorname{Gr}_{\ge0}(k,N)\) une matrice de parcours survivants. Sur le lieu où \(\operatorname{rank}(C_nZ_n^{\mathsf T})=k\), définissons
\begin{equation}
 Y_n=C_nZ_n^{\mathsf T}\in\operatorname{Gr}(k,k+m).
 \label{geometria_positiva:eq:hmt-amplituhedron-map}
\end{equation}

\begin{definition}[Image amplituédrique candidate]
Soit \(\mathfrak C_n^{\rm TPK}\) le complexe positif de parcours et de retenues à profondeur \(n\). Son image
\[
 \mathcal A_n^{\rm HMT}(k,m)
 =\{C Z_n^{\mathsf T}:C\in\mathfrak C_n^{\rm TPK}\}
\]
est appelée image amplituédrique candidate de niveau \(n\). Notons par
\[
 \mathcal F_{{\rm kin},n}:
 \mathfrak C_n^{\rm TPK}\longrightarrow
 \operatorname{Gr}(k,k+m),\qquad
 [C]\longmapsto[CZ_n^{\mathsf T}]
\]
son application cinématique.
\end{definition}

\begin{proposition}[Image directe canonique sous hypothèses de degré et de frontière]
\label{geometria_positiva:thm:hmt-amplituhedron-form} Supposons que \(\mathfrak C_n^{\rm TPK}\) soit un complexe orienté de dimension pure, que \(\mathcal F_{{\rm kin},n}\) soit propre sur sa clôture, génériquement finie de degré un sur chaque composante, préserve l’orientation, couvre son image sans chevauchements intérieurs non déclarés et envoie les frontières sur des frontières de manière compatible avec les résidus. Alors
\begin{equation}
 \Omega_n^{\rm HMT}
 =(\mathcal F_{{\rm kin},n})_*
 \left(\sum_{\sigma\in\mathfrak C_n^{\rm TPK}}\Omega_\sigma\right)
 \label{geometria_positiva:eq:hmt-canonical-pushforward}
\end{equation}
est bien définie, possède des pôles logarithmiques uniquement sur la frontière externe de \(\mathcal A_n^{\rm HMT}\) et est indépendante de toute subdivision ultérieure compatible. Ses résidus sont les formes canoniques des images des faces.
\end{proposition}
\begin{proof}
Le théorème~\ref{geometria_positiva:thm:tpk-higher-residues} annule les facettes intérieures avant l’image directe. Le degré un et l’absence de chevauchements intérieurs empêchent les multiplicités parasites ; la compatibilité frontière--résidu transporte les pôles logarithmiques et leurs résidus. Un raffinement compatible change la décomposition cellulaire, mais non sa somme.
\end{proof}

La proposition détermine une forme structurelle et récupérable uniquement lorsque ses hypothèses sont vérifiées. Si l’on démontre en outre que \(\mathfrak C_n^{\rm TPK}=\operatorname{Gr}_{\ge0}(k,N)\) sur le domaine pertinent et que l’application possède la couverture requise, alors \(\mathcal A_n^{\rm HMT}(k,m)\) coïncide avec l’amplituèdre standard de ces données et \eqref{geometria_positiva:eq:hmt-canonical-pushforward} est sa forme canonique complète. Sans dérivation du complexe, démonstration des hypothèses d’image directe et couverture, l’objet demeure une formalisation candidate ; ce n’est ni un théorème attribué au corpus ni une assertion de surjectivité.

\subsection{Synthèse}

La multimorphologie structurelle obtenue peut s’écrire comme la composition
\[
 \boxed{
 \begin{gathered}\text{histoire TPK}\longmapsto\text{cellule ordonnée}\\
 \longmapsto\text{réalisation positroïdale ordonnée}\\
 \longmapsto\text{forme canonique}\longmapsto\text{archive orientée par }K.\end{gathered}}
\]
Chaque flèche conserve une information différente : le TPK conserve la généalogie ; la grassmannienne conserve les mineurs ; la forme canonique conserve la frontière après annulation des parois internes ; \(K\) conserve le repère métrique et la récupération. La structure discrète du continu apparaît ainsi comme un système admettant simultanément un raffinement infini, une positivité finie à chaque niveau et une limite archivée sans perte de la trajectoire qui l’a engendrée. La réalisation positroïdale, les cartes d’amas, l’extension cubique et l’image amplituédrique sont des formalisations ajoutées avec les statuts précis déclarés dans chaque partie ; elles ne sont pas rétroprojetées comme des théorèmes déjà démontrés par le corpus.

% END OWNER


\section{Stabilité par raffinement et compatibilité des lecteurs sectoriels}

\begin{theorem}[Unité et répartition de l’apport]
\label{geometria_positiva:thm:unidad-distribucion} Le théorème de stabilité par raffinement trouve sa place principale dans la construction conjointe du continu. Ses applications à la forme de Weil, aux cinq réalisations terminales, au couplage arithmético--géométrique et à l’opérateur de masse sont des corollaires typés. Aucune de ces applications ne peut sélectionner rétrospectivement le raffinement ni remplacer sa démonstration.
\end{theorem}
\begin{proof}
La forme canonique et son annulation sont définies avant l’application de tout lecteur sectoriel. Chaque lecteur reçoit la même somme cellulaire et doit commuter avec la frontière et le résidu. La propriété commune est démontrée une seule fois ; la particularité de chaque secteur réside dans son application de réalisation. Si un lecteur sélectionnait les cellules à partir de son résultat, il inverserait la chaîne causale.
\end{proof}

\section{Consecuencias estructurales}

\subsection{Multimorphologie de la grassmannienne et de l’amplituèdre}

Le terme \emph{multimorphologie} désigne ici une structure mathématique précise. Soit \(\widehat{\mathcal X}_{\rm TPK}\) l’espace des états enrichis et soit \(p_{n+1,n}\) son système de troncatures. Une donnée multimorphologique positive est un système de réalisations typées
\[
\begin{aligned}
\mathfrak P_n&:\widehat{\mathcal X}_{\rm TPK,n}
 \longrightarrow \operatorname{Gr}_{\geq0}(k,N_n),\\
\mathfrak R_n&:\widehat{\mathcal X}_{\rm TPK,n}
 \longrightarrow C^{+}_{d,n},\\
\mathfrak C_n&:\widehat{\mathcal X}_{\rm TPK,n}
 \longrightarrow \mathcal Q^{+}_{\rm con,n},\\
\mathfrak M_n&:\widehat{\mathcal X}_{\rm TPK,n}
 \longrightarrow \mathcal L^{+}_{M,n},
\end{aligned}
\]
où les codomaines représentent respectivement la positivité de Plücker, la positivité réticulaire, constitutive et opératoire de masse. Les quatre applications doivent commuter avec leurs liaisons de raffinement et conserver les coordonnées de l’état utilisées par chaque lecture. On n’exige ni l’égalité des quatre codomaines ni la sélection des autres par l’un d’entre eux.

L’application cinématique
\[
\mathfrak A_{Z,n}:C\longmapsto CZ^{\mathsf T}
\]
se compose avec \(\mathfrak P_n\) et produit la morphologie amplituédrique. Ainsi, la grassmannienne organise les cartes, les strates et leurs mutations ; l’amplituèdre organise les images cinématiques et leurs formes canoniques. La multimorphologie est le diagramme maintenant simultanément ces réalisations comme images du même état enrichi.

\begin{proposition}[Cohérence multimorphologique par raffinement]
\label{geometria_positiva:prop:coherencia-multimorfologica} Si chaque réalisation commute avec la troncature et si son application sur les formes commute avec la frontière et le résidu, l’annulation des faces intérieures se transporte aux quatre codomaines et à toute image amplituédrique de degré un. Les formes exprimées sont indépendantes de la profondeur, bien que leurs archives conservent les parcours raffinés.
\end{proposition}
\begin{proof}
La forme de l’état source est la somme orientée de ses cellules. La commutation avec la troncature identifie les deux représentations de chaque raffinement ; la commutation avec la frontière et le résidu transporte l’annulation des incidences opposées. Une application cinématique de degré un préserve la multiplicité de la forme dans chaque intérieur. Seules les frontières externes subsistent donc dans chaque réalisation.
\end{proof}

L’existence de \(\mathfrak R_n\) et de \(\mathfrak M_n\) repose sur les constructions réticulaire et matérielle de référence déjà obtenues. La réalisation \(\mathfrak P_n\) par courbe des moments est exacte comme formalisation ordonnée ; sa promotion en atlas positroïdal intrinsèque exige la mesure de frontière. La réalisation amplituédrique exige en outre le rang, le degré et la couverture. C’est la frontière démonstrative de la multimorphologie, non une restriction de sa définition.

\subsection{Continu}

Le passage à la limite n’ajoute pas de points déconnectés : il stabilise des formes compatibles d’une même structure raffinée. Les parois internes disparaissent de la forme globale, mais demeurent dans la mémoire des trajectoires. On obtient ainsi une distinction exacte entre invariance géométrique et conservation généalogique.

\subsection{Riemann--Weil}

Le télescopage rend la réalisation de la forme arithmétique indépendante du maillage nonadique. La positivité de Weil continue de dépendre des canaux gamma, premier et polaire, ainsi que de l’identité de moments mixtes les transformant en forme de Gram. Le nouveau résultat ne remplace pas cette démonstration : il démontre qu’aucun pôle intérieur du raffinement ne peut introduire une direction négative parasite.

\subsection{Cinq réalisations terminales}

Un lecteur terminal transporte l’annulation s’il commute avec la frontière et si le résidu de la forme transportée coïncide avec la forme de la face. Cette condition unifie le mécanisme de recollement sans dissocier le continu ni transformer les cinq problèmes en cinq générateurs.

\subsection{Couplage et flot}

La constante holographique de couplage arithmético--géométrique récupère \(K\), et \(K\) détermine \(O_K\), \(G_K\) et l’archive. L’action générale \(D_K(s)\) ouvre une voie positive. La direction centrée \(u_K=P_3P_{11}K\), satisfaisant \(u_K\perp\mathbf1\), permet de l’affiner en
\[
g_{K,\rho}(\tau)=
\operatorname{diag}
\bigl(\rho_A^{\tau\widehat u_{K,1}},\ldots,
\rho_A^{\tau\widehat u_{K,12}}\bigr)\in SL(12,\mathbb R)_{>0}.
\]
Les rapports \(\varphi^{-2}\), \(\pi/\varphi^2\) et \(\varphi/\pi^2\) peuvent paramétrer différentes sections du flot après avoir été engendrés et typés comme expressions de la quadrirelation coinductive \((\pi,\varphi,e,\alpha)\) ; aucun ne sélectionne rétrospectivement le TPK ni ne se confond avec \(\kappa_{\rm ag}\).

\subsection{Particules et masses}

La congruence
\[
M_M^{\beta,r}
=R_{\rm act}^{B_M^r/2}M_M^{(0)}R_{\rm act}^{B_M^r/2}
\]
préserve le cône positif. L’exception électronique est antérieure à cette évaluation : elle provient de l’unique point fixe \((5,5)\). La frontière \(111111\) et la signature \(555555\) demeurent des registres distincts de la même généalogie.

\section{Critère de promotion en atlas positroïdal et en amplituèdre}

La promotion sera achevée lorsque seront concrètement construits :
\begin{enumerate}
\item un réseau planaire \(\mathcal N_z\) par cylindre, avec des poids positifs dérivés du TPK ;
\item la matrice de mesure de frontière et sa formule LGV pour tous les mineurs ;
\item la permutation décorée ou le collier de Grassmann de chaque strate ;
\item une graine ou un graphe plabique dont la mutation commute avec le raffinement et la mémoire ;
\item l’application cinématique \(C\mapsto CZ^{\mathsf T}\), avec rang, dimension, degré, couverture et correspondance des frontières ;
\item l’identité d’image directe des formes canoniques.
\end{enumerate}
Ces conditions ne constituent pas une liste sociologique : ce sont les critères de réfutation distinguant une réalisation complète d’une collection de points positifs.

\section{Conclusion}

L’apport central est une loi de stabilité géométrique du raffinement : la profondeur ajoute de la mémoire et de la résolution, tandis que la forme canonique conserve la frontière totale. Le registre \(K\) rend ce processus récupérable et métrique ; sa direction centrée organise un flot de déterminant un respectant les signes de Plücker. L’électron occupe l’unique centre équivariant de la branche matérielle. Avec ces éléments, le passage à la grassmannienne positive cesse d’être une analogie : il possède une réalisation ordonnée exacte et un programme de promotion dont les applications et les critères de réfutation sont entièrement spécifiés.

